\chapter{Postmenopausal Bleeding}

\section*{Definition}
Postmenopausal bleeding (PMB) is any vaginal bleeding occurring 12 or more months after the final menstrual period. It must be considered endometrial cancer until proven otherwise, although 90--95\% of cases are due to benign causes.

\section*{What Happens in Your Body}
The endometrium atrophies after menopause due to estrogen deprivation, leaving fragile vessels prone to spontaneous rupture. When endometrial proliferation occurs (from exogenous estrogen, endogenous estrogen production via peripheral aromatization of androstenedione, or estrogen-secreting tumors), the thickened endometrium outgrows its blood supply and sheds irregularly.

\section*{Causes}
\begin{itemize}
  \item \textbf{Benign uterine:} Endometrial or cervical atrophy (most common), endometrial polyps, submucosal fibroids, endometritis
  \item \textbf{Malignant/premalignant:} Endometrial hyperplasia (with or without atypia), endometrial adenocarcinoma (most common gynecologic malignancy), uterine sarcoma, cervical carcinoma, vaginal carcinoma
  \item \textbf{Hormonal:} Unopposed estrogen therapy (without progestin), tamoxifen (partial estrogen agonist on endometrium), phytoestrogen supplements, endogenous estrogen from obesity or granulosa cell tumors
  \item \textbf{Miscellaneous:} Vaginal atrophy, trauma (including post-coital bleeding), anticoagulant use, lichen sclerosus, cervical polyps
\end{itemize}

\section*{When to See a Doctor}
See your doctor if:
\begin{itemize}
  \item Any vaginal bleeding after 12 months of amenorrhea (all episodes require evaluation)
  \item Spotting, pink-tinged discharge, or blood on toilet paper in a postmenopausal woman
  \item Post-coital bleeding in a postmenopausal person
  \item Concurrent abdominal or pelvic pain, bloating, or unexplained weight loss
\end{itemize}

\section*{Related Symptoms}
\begin{itemize}
  \item Vaginal spotting, pelvic pressure or pain, dyspareunia, abdominal bloating, unexpected weight loss, fatigue (if anemia develops), purulent or bloody vaginal discharge
\end{itemize}
\textbf{Important:} The risk of endometrial cancer is approximately 5--10\% in women with postmenopausal bleeding; this risk increases with age, obesity, diabetes, hypertension, and prolonged unopposed estrogen exposure.

\section*{Self-Care / Home Management}
\begin{itemize}
  \item Maintain a bleeding diary recording date, duration, and volume (pad/tampon count)
  \item Avoid vaginal intercourse, tampons, and douching until evaluation is complete
  \item Use panty liners for hygiene; do not use any vaginal estrogen preparations until malignancy is excluded
  \item Continue prescribed medications (especially anticoagulants) unless directed otherwise
\end{itemize}

\section*{How Your Doctor Will Evaluate You}
\begin{itemize}
  \item Transvaginal ultrasound (TVUS) to measure endometrial thickness (ET): ET \textless{}4--5 mm has \textgreater{}99\% negative predictive value for endometrial cancer; ET \textgreater{}5 mm requires tissue sampling
  \item Endometrial biopsy (office Pipelle): first-line tissue diagnosis with 90--98\% sensitivity for endometrial cancer
  \item Hysteroscopy with directed biopsy if office biopsy is insufficient, nondiagnostic, or if focal lesions (polyps, submucosal fibroids) are suspected
  \item Saline infusion sonohysterography for further characterization of intracavitary lesions
  \item Laboratory: CBC (evaluate for anemia), coagulation panel if menorrhagia-like pattern
\end{itemize}

\section*{Treatment Options}
\begin{itemize}
  \item \textbf{Atrophic endometrium:} Reassurance; consider vaginal moisturizers or low-dose vaginal estrogen after malignancy is excluded
  \item \textbf{Endometrial polyp/fibroid:} Hysteroscopic polypectomy or myomectomy
  \item \textbf{Endometrial hyperplasia without atypia:} Medroxyprogesterone acetate or levonorgestrel-releasing IUD (Mirena) with repeat biopsy in 3--6 months
  \item \textbf{Endometrial hyperplasia with atypia / endometrial cancer:} Total hysterectomy with bilateral salpingo-oophorectomy; staging with lymph node dissection as indicated; referral to gynecologic oncology
  \item \textbf{Hormonal manipulation:} Discontinue or adjust exogenous estrogen; switch tamoxifen to an aromatase inhibitor if feasible for breast cancer people
\end{itemize}

\section*{References}
\begin{thebibliography}{99}
\bibitem{postmenopausal_bleeding:ref1} Clarke MA, Long BJ, Del Mar Morillo A, et al. ``Endometrial carcinoma risk among women with postmenopausal bleeding: a systematic review and meta-analysis.'' \textit{JAMA Internal Medicine}, 2018;178(9):1210--1222.
\bibitem{postmenopausal_bleeding:ref2} Smith-Bindman R, Weiss E, Feldstein V. ``How thick is too thick? When endometrial thickness should prompt biopsy in postmenopausal women without vaginal bleeding.'' \textit{Ultrasound in Obstetrics \& Gynecology}, 2004;24(5):558--565.
\bibitem{postmenopausal_bleeding:ref3} ACOG Committee Opinion No. 734. ``The role of transvaginal ultrasonography in evaluating the endometrium of women with postmenopausal bleeding.'' \textit{Obstetrics \& Gynecology}, 2018;131(5):e124--e129.
\bibitem{postmenopausal_bleeding:ref4} Goldstein RB, Bree RL, Benson CB, et al. ``Evaluation of the woman with postmenopausal bleeding: Society of Radiologists in Ultrasound-Sponsored Consensus Conference statement.'' \textit{Journal of Ultrasound in Medicine}, 2001;20(10):1025--1036.
\bibitem{postmenopausal_bleeding:ref5} Bakkum-Gamez JN, Gonzalez-Bosquet J, Laack NN, et al. ``Current issues in the management of endometrial cancer.'' \textit{Mayo Clinic Proceedings}, 2008;83(1):97--112.
\bibitem{postmenopausal_bleeding:ref6} Gupta JK, Chien PFW, Voit D, et al. ``Ultrasonographic endometrial thickness for diagnosing endometrial pathology in women with postmenopausal bleeding: a meta-analysis.'' \textit{Acta Obstetricia et Gynecologica Scandinavica}, 2002;81(9):799--808.

\end{thebibliography}
